\documentclass[10pt,a4paper]{amsart}
\usepackage[margin=19mm]{geometry}
\usepackage{amsmath,amssymb,mathtools}
\usepackage[scr=rsfs,cal=euler]{mathalfa}
\usepackage{xfrac}
\usepackage[unicode,hidelinks,bookmarks=false]{hyperref}
\newcommand{\ie}{i.e.\ }
\newcommand{\cf}{cf.\ }
\newcommand{\resp}{respectively\ }
\newcommand{\Subsection}{\subsection}
\providecommand{\nobreakdash}{\nobreak\hbox{-}}
\setlength{\emergencystretch}{2em}
\multlinegap0pt
\usepackage{bm}
\usepackage{bbm}
\usepackage{dsfont}
\usepackage{xspace}
\usepackage{cancel}
\usepackage{mathtools}
\usepackage{amsthm}
\makeatletter
\@ifundefined{theorem}{\newtheorem{theorem}{Theorem}}{}
\@ifundefined{lemma}{\newtheorem{lemma}[theorem]{Lemma}}{}
\makeatother
\title[Invariant part of class group and distribution of relative c]{Invariant part of class group and distribution of relative class group}
\author{Weitong Wang}
\date{Source: arXiv:2408.06676v2 (CC BY 4.0)}
\begin{document}
\maketitle

\noindent\emph{Source and licence.} Invariant part of class group and distribution of relative class group. Version arXiv:2408.06676v2 (CC BY 4.0), distributed under the Creative Commons Attribution 4.0 International licence, as recorded in the arXiv licence metadata for this exact version and shown on the abstract page https://arxiv.org/abs/2408.06676v2. Full rights notice: licenses/arxiv-2408.06676-CC-BY-4.0.txt.\par\smallskip

\noindent\textit{Editorial scope.} Counted unit: the Lemma whose source label is \texttt{lemma of the lemma} in the article (source0.tex lines 2390--2410, source label \texttt{lemma of the lemma}), delivered together with the article's own complete original proof of it (source0.tex lines 2411--2477, one \texttt{proof} environment of 67 lines). No mathematics is written, repaired or re-derived: statement and proof are the article's own text line for line. Required context retained uncounted: none. Every definition, display and label of those lines is retained unchanged. Omitted from this package and not redistributed: 1--2389, 2478--3124. Nothing inside a retained excerpt is omitted or altered: every retained body is delivered exactly as the article's lines are written, so no omission receipt of any kind accompanies this package. Citations: the retained excerpts carry no citation command at all, so no bibliography entry of the article is transcribed and no reference is redistributed. Reference adaptation: none was needed and none was made; every reference inside the delivered unit resolves inside it, so no unresolved reference or citation is produced. Printed slips kept literal: no printed word, symbol, hypothesis or number of the counted statement or of its proof was altered. Layout and class: the article's own class is \texttt{amsart} with packages including amsmath, amssymb, amscd, verbatim, xspace, amsthm, tensor, graphicx, enumerate, tikz, tikz-cd, appendix, latexsym, epsfig; the wrapper is the lane's pinned \texttt{amsart} editorial wrapper at 10pt on A4, which restates the theorem-like environments the retained text needs and adds the article's own macro definitions that the retained text needs (none), with extra wrapper packages (bm, bbm, dsfont, xspace, cancel, mathtools). The delivered numbering is the unit's own. Assets: no retained span contains a figure environment, a graphics-inclusion command, a PDF-page inclusion, a table or a TikZ picture; no image file is copied into this package, the package declares no asset dependency and no image file is redistributed with it. Licence: version arXiv:2408.06676v2 of this article is distributed under the Creative Commons Attribution 4.0 International licence, as recorded in the arXiv licence metadata for this exact version and shown on the abstract page https://arxiv.org/abs/2408.06676v2. A full rights notice is delivered at licenses/arxiv-2408.06676-CC-BY-4.0.txt. Characters: every retained excerpt is pure ASCII, so the delivered engine, XeLaTeX with its default Latin Modern text font, reports no missing character.
\begin{lemma}\label{lemma of the lemma}
	\begin{enumerate}
		\item Fix an abelian $\bar{G}$-field $K/\mathbb{Q}$, and define $\mathcal{E}_K:=\{L\in\mathcal{E}\mid K\subseteq L\}$.
		We have
		\begin{equation*}
			N_C(\mathcal{E}_K,X)\leq\lvert\operatorname{Hom}_{\bar{G}}(\operatorname{Cl}_K,H)\rvert\cdot\#\{\chi\in\operatorname{Hom}_{\bar{G}}(J_K^{S_{\infty}}/K^{S_{\infty}}, H)\mid C(\chi)<X\},
		\end{equation*}
		where $S_{\infty}$ is the set of all infinite primes of $K$.
		\item Let $A$ and $B$ be two finite abelian $p$-groups.
		We have
		\begin{equation*}
			\lvert \operatorname{Hom}(A,B)\rvert\leq\lvert A\rvert^{\operatorname{rk}_p B}.
		\end{equation*}
		\item Fix a finite discrete $\bar{G}$-module $A$.
		There exist positive constant $c$ and $N$ depending only on $\bar{G}$ and $A$ such that for each abelian $\bar{G}$-field $K$, we have
		\begin{equation}\label{eqn: lemma of the lemma 1}
			\lvert\operatorname{Hom}(\operatorname{Cl}_K,A)\rvert\leq c\sqrt{d_K}^{N},
		\end{equation}
		where $\sqrt{d_K}$ is the radical of the absolute discriminant (product of ramified primes).
	\end{enumerate}
\end{lemma}

\begin{proof}
	(1):
	Recall that $S_{\infty}$ is the set of infinite primes of $K$.
	For a finite set $S$ of primes including the ones at infinity, we have the notions of $S$-id{\`e}les $J_K^{S}=\prod_{\mathfrak{p}\notin S}U_{\mathfrak{p}}\times\prod_{\mathfrak{p}\in S}K_{\mathfrak{p}}^*$, and $S$-units $K^{S}=K^*\cap J_K^{S_\infty}$ (inside $J_K$).
	In addition, there is a short exact sequence
	\begin{equation*}
		1\to J_K^{S}/K^S\to\operatorname{C}_K\to\operatorname{Cl}_K^S\to1.
	\end{equation*}
	When $S=S_{\infty}$, we have $\operatorname{Cl}_K^{S_\infty}=\operatorname{Cl}_K$.
	By taking the homomorphisms that are $\bar{G}$-equivariant to $H$, we have an exact sequence
	\begin{equation*}
		0\to\operatorname{Hom}_{\bar{G}}(\operatorname{Cl}_K,H)\to\operatorname{Hom}_{\bar{G}}(\operatorname{C}_K,H)\to\operatorname{Hom}_{\bar{G}}(J_K^{S_{\infty}}/K^{S_{\infty}},H).
	\end{equation*}
	This implies first that if $\chi:J_K^{S_{\infty}}/K^{S_{\infty}}\to H$ admits a preimage, then the number of $\tilde{\chi}:\operatorname{C}_K\to H$ such that its restriction is exactly $\chi$ is given by the size of $\operatorname{Hom}_{\bar{G}}(\operatorname{Cl}_K,H)$.
	In particular, this shows that the original proof of Lemma 7.10 is flawed.
	That is, when $\operatorname{Hom}_{\bar{G}}(\operatorname{Cl}_K,H)$ is non-trivial, there must be multiple $G$-fields $L/\mathbb{Q}$ including $K$ as a subfield with exactly the same local specifications.
	
	On the other hand, we have shown that $\mathcal{E}_K$ and $\operatorname{Sur}_{\bar{G}}(\operatorname{C}_K,H)$ corresponds to each other bijectively (see Proposition 7.6).
	And we see in Section 7.1 that the counting function $C$ could be generalized to $\chi:J_K^{S_{\infty}}\to H$ (that factors through $K^{S_{\infty}})$.
	Putting everything together, we have
	\begin{equation*}
		\begin{aligned}
			N_C(\mathcal{E}_K,X)=&\#\{\tilde{\chi}\in\operatorname{Sur}_{\bar{G}}(\operatorname{C}_K,H)\mid C(\tilde{\chi})<X\}
			\\
			\leq&\lvert\operatorname{Hom}_{\bar{G}}(\operatorname{Cl}_K,H)\rvert\cdot\#\{\chi\in\operatorname{Hom}_{\bar{G}}(J_K^{S_{\infty}}/K^{S_{\infty}}, H)\mid C(\chi)<X\}.
		\end{aligned}
	\end{equation*}
	And we are done for (1).
	
	(2):
	For finite cyclic $p$-groups, we see that
	\begin{equation*}
		\operatorname{Hom}(\mathbb{Z}/p^a\mathbb{Z},\mathbb{Z}/p^b\mathbb{Z})\cong\mathbb{Z}/p^{\min\{a,b\}}\mathbb{Z}.
	\end{equation*}
	In general, write $A\cong\prod_{i=1}^{\operatorname{rk}_p A}\mathbb{Z}/p^{a_i}\mathbb{Z}$ and $B\cong\prod_{i=1}^{\operatorname{rk}_p B}\mathbb{Z}/p^{b_i}\mathbb{Z}$, we have
	\begin{equation*}
		\operatorname{Hom}(A,B)\cong\prod_{i=1}^{\operatorname{rk}_p A}\prod_{j=1}^{\operatorname{rk}_p B}\mathbb{Z}/p^{\min\{a_i,b_j\}}\mathbb{Z}.
	\end{equation*}
	This implies that
	\begin{equation*}
		\lvert\operatorname{Hom}(A,B)\rvert\leq\prod_{j=1}^{\operatorname{rk}_p B}\prod_{i=1}^{\operatorname{rk}_p A}p^{a_i}=\lvert A\rvert^{\operatorname{rk}_p B}.
	\end{equation*}
	And we are done for (2).
	
	(3):
	By (2), we have
	\begin{equation*}
		\begin{aligned}
			\lvert\operatorname{Hom}_{\bar{G}}(\operatorname{Cl}_K,A)\rvert\leq&\lvert\operatorname{Hom}(\operatorname{Cl}_K,A)\rvert\\
			=&\prod_{p}\lvert\operatorname{Hom}(\operatorname{Cl}_K[p^\infty],A[p^\infty])\rvert
			\\
			\leq&\prod_p\lvert\operatorname{Cl}_K[p^\infty]\rvert^{\operatorname{rk}_p A}
			\leq\lvert\operatorname{Cl}_K\rvert^{r_A},
		\end{aligned}
	\end{equation*}
	where $r_A=\sup_p\operatorname{rk}_pA$.
	By Minkowski bound of the size of the class group, for each $\varepsilon>0$, there exists a constant $c=c(\lvert \bar{G}\rvert,\varepsilon)$ such that
	\begin{equation*}
		\lvert\operatorname{Cl}_K\rvert^{r_A}\leq c d_K^{\frac{1+\varepsilon}{2}r_A},
	\end{equation*}
	where $d_K$ is the absolute discriminant of $K/\mathbb{Q}$.
	Since the Galois group $\bar{G}$ is fixed, we see that there exists $N=N(\bar{G},\varepsilon,A)$ such that $(\sqrt{d_K})^N>d_K^{\frac{1+\varepsilon}{2}r_A}$, where $\sqrt{d_K}$ means the radical of the absolute discriminant (or just the product of ramified primes), that is, 
	\begin{equation*}
		\lvert\operatorname{Hom}(\operatorname{Cl}_K,A)\rvert\leq c\sqrt{d_K}^{N},
	\end{equation*}
	and we are done for the proof.
\end{proof}

\end{document}
